\documentclass[11pt]{article}
\usepackage[margin=1.15in]{geometry}
\usepackage{amsmath,amssymb,amsthm}
\usepackage[expansion=false]{microtype}
\usepackage[hidelinks]{hyperref}
\usepackage{setspace}
\onehalfspacing

\newtheorem{proposition}{Proposition}
\theoremstyle{definition}
\newtheorem{definition}{Definition}

\title{Decadence as Transformation:\\Nietzsche, Voltaire, and Rousseau on Decline}
\author{Flyxion\\\small Independent Researcher}
\date{September 2026}

\begin{document}
\maketitle

\begin{abstract}
\noindent Nietzsche is often read as a theorist of civilizational fall, a thinker who mourns a lost health and blames modern society for its loss. The reading fails at exactly the point where he sets Voltaire against Rousseau. This essay argues that Nietzsche's vocabulary of decadence, sickness, and degeneration operates inside a framework with no conserved original state, so that decline means a change in the organization of forces and never an expulsion from innocence. Three claims are developed. First, Rousseau's genealogy and Nietzsche's differ in structure: one has an invariant object that is corrupted, the other has transformation throughout. Second, Nietzsche reads Rousseau's doctrine of natural goodness as the conversion of a grievance into a universal principle, an operation that gives resentment the standing of insight. Third, Nietzsche treats doctrines as symptoms, so that the appeal of Rousseau to the suffering is itself evidence to be interpreted. An appendix applies the internal criterion to contemporary cases (protective upbringing, entertainment, religion) as extensions and not as claims about Nietzsche's own text. The essay closes by marking where the Nietzschean reading is partial, both as an account of Rousseau and as a self-description, since Nietzsche's own diagnosis can be read as a symptom by the method he recommends.
\end{abstract}

\section{The Misplaced Fall}

A reader who meets Nietzsche's language of exhaustion, herd morality, and the last man may reasonably conclude that he holds a narrative of decline with a recognizable shape: something was once whole, something intervened, and the present is the wreckage. That shape is the shape of Rousseau's story, and Nietzsche's contrast between Rousseau and Voltaire exists in large part to refuse it.

The refusal is not a refusal to speak of decline. Nietzsche says that cultures weaken, that valuations turn against the conditions that produced them, and that some forms of life become exhausted. What he denies is the background assumption that makes decline a fall. There is no innocent term at the start of the sequence. There are drives, cruelty, appropriation, hierarchy, imitation, fear, and creative aggression, and what civilization does with them is organize, redirect, suppress, sublimate, and occasionally intensify.\footnote{Cf.\ \emph{Beyond Good and Evil}, \S\S13, 23, and 36 on strength, drives, and Nietzsche's proposed drive-psychology; \S259 on appropriation, incorporation, suppression, and exploitation as characteristics of life; and \S229 on the spiritualization and deepening of cruelty on which higher culture rests. On cruelty in the origins of guilt and debt, cf.\ \emph{On the Genealogy of Morality}, Second Essay, \S6.} The question is never how a good creature was ruined. It is what configuration a given organization of forces produces, and what that configuration permits.

The dedication of \emph{Human, All Too Human} to the memory of Voltaire marks the methodological allegiance.\footnote{The dedication appears only in the first edition of 1878, marking the centenary of Voltaire's death; it was dropped when Nietzsche reissued the work in two volumes in 1886. References to Nietzsche's works in this essay use the section numbers shared by the standard English translations (Kaufmann, Hollingdale, Polt); the wording of individual phrases varies by translation.} Voltaire supplies a temperament: skeptical, observant, hostile to fanaticism, unwilling to purchase consolation with metaphysics. Rousseau supplies the temptation Nietzsche most distrusts, which is the consolation of an origin.

\section{Two Genealogies}

The difference between the two accounts, as Nietzsche constructs it, can be stated structurally. Rousseau's account, so construed, contains a conserved object. Let $H_{\mathrm{nat}}$ denote natural humanity and $H_{\mathrm{cor}}$ corrupted humanity. Then
\[
H_{\mathrm{nat}} \xrightarrow{\ \text{society}\ } H_{\mathrm{cor}},
\]
and critique is the attempt to recover $H_{\mathrm{nat}}$, or at least to measure the present against it. The origin functions as a standard. Whatever deviates from it is degraded by definition.

Nietzsche's account removes the standard. Let $D=\{d_1,\dots,d_n\}$ be a population of drives, and let a social order act on it by a transformation $T$ that inhibits, ranks, redirects, sublimates, or internalizes:
\[
D_t \xrightarrow{\ T_t\ } D_{t+1} \xrightarrow{\ T_{t+1}\ } D_{t+2} \cdots
\]
No privileged $D_0$ is designated as the healthy baseline to which the sequence should return. The history of a people is a chain of configurations, each produced from the last by particular mechanisms.\footnote{The textual warrant is \emph{On the Genealogy of Morality}, Second Essay, \S12, where Nietzsche separates the origin of a thing from its later purpose, since existing things are repeatedly appropriated and reinterpreted by new powers, and adds that partial loss of utility and even degeneration can belong to genuine progress.}

\begin{definition}[Conserved-origin genealogy]
A genealogy is \emph{conserved-origin} if it names a term $H_0$ such that later terms are evaluated by their distance from $H_0$ and critique aims at restoring or approximating $H_0$.
\end{definition}

\begin{definition}[Transformational genealogy]
A genealogy is \emph{transformational} if every term is the product of a prior transformation, no term is privileged as a norm by its priority, and evaluation attaches to configurations by what forms of life they permit.
\end{definition}

Rousseau's story, as Nietzsche receives it, is of the first kind. Nietzsche's is of the second. The consequence is that the vocabulary of decline changes meaning. In a conserved-origin genealogy, decline is distance from the origin. In a transformational one, decline must be defined internally, by reference to what a configuration does to its own conditions of existence.

\section{Decadence as Internal Failure}

The internal definition is the interesting one. A configuration of valuations can be generative, in that it sustains the drives and tensions that produced it and continues to yield demanding forms of life. It can also become self-negating, rewarding safety, conformity, guilt, universalized pity, or the suppression of exceptional forms, until the tensions that made anything difficult possible have been bled away.

\begin{proposition}
On the transformational account, decadence is a relation between a configuration and its own conditions of production, not between a configuration and an origin.
\end{proposition}

This is why Nietzsche can admire civilization and attack it in the same breath without inconsistency. Enormous constraint and sublimation are required to produce extraordinary cultural forms.\footnote{The clearest text is \emph{On the Genealogy of Morality}, Second Essay, \S\S16--18, where instincts that can no longer discharge outward are turned inward, producing bad conscience, which Nietzsche calls a grave illness and also the womb of ideal and imaginative events. Goethe in \emph{Twilight of the Idols}, ``Skirmishes,'' \S49, is the corresponding positive case, a type who takes opposing forces into himself and disciplines them into a whole.} The same machinery can also produce domestication. Civilization is neither progressive nor corrupting as such. It is a transformation, and the question of value attaches to the output of a given transformation and to whether that output can sustain what generated it.

The same structure explains his rejection of the opposite story. If the Rousseauian claim is that we fell from something good, the Enlightenment-progressive claim is that we are advancing toward something better by the mere accumulation of reason and comfort. Nietzsche rejects both, for the same reason: each assumes a direction built into history. A transformational account offers no such guarantee. Some configurations open new possibilities and some close them.

A consequence follows that the essay's formal apparatus, introduced in the conclusion, makes precise. Because decadence is defined by the relation of a configuration to its own conditions of production, it is a property of a trajectory and not of a state. Two societies could show identical present levels of comfort, constraint, artistic production, religious observance, and leisure while one regenerates its capacities and the other depletes them. The snapshot does not settle the diagnosis. History and direction are required, so no detector of decadence can be built from surface features alone. Luxury is not decadence, pleasure is not decadence, protection is not decadence, and by the same token constraint is not health and hardship is not health. Decadence is not a state of a system but a sign of its derivative (see equation~\eqref{eq:delta}).

\begin{proposition}[Comparative constraint]
A diagnosis of decadence cannot be inferred from the presence of comfort, protection, pleasure, constraint, or suffering alone. It requires evidence that, relative to an appropriate comparator, the configuration persistently reduces the regeneration of capacities relevant to the diagnosis.
\end{proposition}

The proposition keeps the internal criterion from collapsing into taste. Without comparative evidence, ``decadence'' names the evaluator's valuation of a configuration. With such evidence it can also name a claim about the configuration's dynamics. A negative diagnosis therefore cannot rest on the evaluator's say-so. It needs an observable transition, such as a documented decline in what the configuration can still do, that can be checked independently of the judgment it supports, so that the alleged loss of regenerative capacity is distinguishable from mere disapproval of the state.

\section{From Grievance to Principle}

Nietzsche's hostility to Rousseau is not only theoretical. In the section of \emph{Twilight of the Idols} titled ``Skirmishes of an Untimely Man,'' he attacks Rousseau in connection with the moral atmosphere of the Revolution and the doctrine of equality.\footnote{\emph{Twilight of the Idols}, ``Skirmishes of an Untimely Man,'' \S48, which calls Rousseau idealist and rabble in one person and names the doctrine of equality the most poisonous of poisons. The Rousseau passage follows Nietzsche's remark on Napoleon in the same section.} The pairing with Voltaire is a structural device of this essay, anchored in the dedication discussed in Section~1, and not a single passage in which Nietzsche sets the two against each other. The charge that Rousseau was a kind of populist agitator is a fair short summary of the political register, provided it is understood structurally.

The structural claim concerns an operation on grievance. Consider the sequence
\[
\text{subordination}
\rightarrow
\text{resentment}
\rightarrow
\text{moral interpretation}
\rightarrow
\text{universal principle}
\rightarrow
\text{political legitimacy}.
\]
The decisive step is the third. The statement ``those above me have advantages I lack'' is politically weak and psychologically uncomfortable, since it concedes the standing of what it resents. The statement ``their superiority is artificial and our equality is natural'' is a different kind of sentence. The particular interest has vanished into a claim about humanity. Resentment is now insight, and the lower position in a hierarchy becomes presumptive evidence of innocence, while refinement, rank, and cultivation become presumptive evidence of corruption.

The pair of oppositions
\[
\text{natural} / \text{ordinary} / \text{uncorrupted}
\quad\text{versus}\quad
\text{artificial} / \text{elite} / \text{decadent}
\]
carries a second message under the first. The apparent description of human nature also allocates moral credit. This is the sense in which a category that looks descriptive can operate as a normative weapon. The complaint is therefore stronger than the observation that Rousseau appealed to the masses. It is that he helped construct a language in which a constituency could experience its historically situated passions as the voice of humanity itself.

Nietzsche's question in the \emph{Genealogy} runs the same way.\footnote{\emph{On the Genealogy of Morality}, Preface \S3 (under what conditions human beings devised the value judgments good and evil) and Preface \S6 (the value of the values); on resentment becoming creative and giving values, First Essay \S10.} It asks not only whether a proposition $P$ is true but which configuration of forces gains authority when $P$ is accepted, and who acquires the power to judge whom.

\section{The Symptom and the Argument}

The second element is the treatment of doctrines as symptoms. Nietzsche repeatedly reads a philosophy as an expression of the condition of the person, or type, for whom it is compelling.\footnote{Nietzsche supplies the methodological vocabulary himself in \emph{On the Genealogy of Morality}, Preface \S6, where morality is to be investigated as result, symptom, mask, and disease, and equally as cause, remedy, stimulant, fetter, and drug. On the type-psychology and naturalism behind the method, see also Leiter, \emph{Nietzsche on Morality}.} The inquiry then has two questions in place of one: whether the doctrine is true, and why it feels true to those to whom it does.

One can write the second question as a conditional,
\[
P(\text{finds Rousseau compelling}\mid \text{condition of life}),
\]
and Nietzsche's claim is that this quantity is higher for those who experience society as constraining, humiliating, or painful. For them Rousseau supplies an interpretation under which suffering becomes evidence against the existing order and not a fact about the sufferer's situation to be investigated. The story has a consoling structure:
\[
\text{present suffering}
\rightarrow
\text{society is corrupt}
\rightarrow
\text{humanity was, or naturally is, better}
\rightarrow
\text{abolish the corrupting arrangement}.
\]
Each arrow is an inference Nietzsche declines. Suffering under an arrangement does not show that the arrangement is a degeneration from a healthier natural state, and it does not by itself reveal the causal history of the suffering. The language of sickness in this context should be read with care: the claim concerns the fit between a doctrine and a condition, and it does not by itself refute the doctrine.

That last point is a constraint on the method, and it should be stated plainly. Explaining why a belief is attractive to a class of believers is not a demonstration that the belief is false. The genetic fallacy is available to any reader who treats the symptom reading as a refutation. Nietzsche's more careful moves place the symptom reading beside independent argument, and use it to explain why an inference that is not compelling nonetheless persuades.

\section{Voltaire's Method}

Voltaire figures in this picture as the model of emancipation by disillusionment. Remove theological, moral, and metaphysical falsifications, and the mind becomes freer. The operation is subtractive, but it subtracts falsehood and does not subtract civilization. It requires no innocent baseline against which to measure the result. Cruelty, superstition, clerical power, and stupidity can be criticized on grounds that need no lost paradise.

Rousseau's operation looks superficially similar, since it also removes falsifications, but it differs in what it takes to lie beneath them. Stripping away the artifice of society, Rousseau expects a nature to appear. Nietzsche expects nothing so accommodating. What lies beneath the moral vocabulary is a field of forces, and forces have no obligation to be benign.

Nietzsche wants the first operation without the comfortable Enlightenment conclusion that disillusionment must yield improvement. This is the reason the Voltaire and Rousseau pair does not sort into rationalist versus romantic. The relevant division is between a skepticism that punctures moral mythology and a sentiment that can manufacture it.

\section{Revaluation and Restoration}

If the present is diagnosed as decadent without an innocent origin behind it, the remedy cannot be restoration. There is no earlier state to restore, and the attempt to reconstruct one would merely install a new fiction. Restoration also meets a formal obstacle. A transformation $T:X_t\mapsto X_{t+1}$ need not have an inverse. Sublimation and internalization are many-to-one, in that different prior configurations can yield the same later one, and information about the earlier condition is lost in the passage. Where no inverse exists, going back is not an operation that can be performed at all, so that the appeal to a recoverable prior state is undefined and not merely undesirable. What remains is revaluation: the production of a configuration that may never have existed. The pattern is
\[
\text{formation} \rightarrow \text{stabilization} \rightarrow \text{inversion or exhaustion} \rightarrow \text{possible revaluation},
\]
a sequence in which the last term is an invention and not a return. Nietzsche's own phrase for the direction is revealing. In \emph{Twilight of the Idols} he allows himself to speak of a return to nature, and then corrects the phrase: it is not a going back but a climbing up.\footnote{\emph{Twilight of the Idols}, ``Skirmishes of an Untimely Man,'' \S48 (``Progress in my sense''), where Napoleon is offered as an instance of the return to nature so understood.} The correction is the entire argument compressed into a sentence.

\section{Limits of the Reading}

Three qualifications keep the account honest.

First, Nietzsche's Rousseau is a partial Rousseau. In the \emph{Discourse on Inequality},\footnote{Rousseau, \emph{Discourse on the Origin and Foundations of Inequality Among Mankind}, Preface and Exordium (the state of nature as conditional reasoning and not historical fact). For a reading that resists the caricature, see Neuhouser, \emph{Rousseau's Theodicy of Self-Love}.} Rousseau describes the state of nature as a hypothesis that may never have obtained, offered as a method for separating what is original in human beings from what is acquired. He does not propose a literal return, and his later political writings accept that society cannot be undone.\footnote{Rousseau, \emph{Of the Social Contract}, Book I, ch.~1, which takes the passage to civil society as given and asks under what conditions it can be legitimate.} A reader who takes Nietzsche's portrait as an exposition of Rousseau will therefore misdescribe him. The portrait is better read as an account of a tendency that the Rousseauian idiom made available, and that later movements exploited, whatever Rousseau intended.

Second, Nietzsche's own vocabulary sometimes betrays the position. Words such as degeneration and physiological decline carry a biologistic weight that sits uneasily beside the claim that there is no baseline. The internal definition of decadence offered in Section 3 is a charitable reconstruction that removes the tension. It is not always the reading the texts invite.

Third, the symptom method is reflexive. If a doctrine's appeal to the suffering counts against its authority, then Nietzsche's revaluative program can be asked the same question: for whom is the diagnosis of decline compelling, and what does the diagnosis do for them? A reader sympathetic to the method should welcome the question, since the method's whole point is that no position is exempt.

Voltaire, finally, is not simply Nietzsche's ally. His deism and his moderate faith in improvement are part of what Nietzsche sets aside when he appropriates Voltaire's temperament. The Voltaire of the essay is a role in Nietzsche's argument as much as a historical figure.

\section{Conclusion}

The difference between Rousseau and Nietzsche is not adequately described as optimism against pessimism, nature against culture, or equality against hierarchy. It concerns the geometry of historical explanation. In the Rousseau Nietzsche constructs as his opponent, history has a privileged point of comparison. Natural humanity precedes corruption, and the injuries of the present become intelligible as departures from an earlier condition. Nietzsche removes that point. There is no position outside transformation from which transformation itself can be judged.

This removal changes the meaning of decadence. If there is no conserved origin, then decadence cannot mean distance from an original health. It must name a relation internal to a configuration of forces, and because that relation concerns a configuration's conditions of production, it is dynamic. Three modes of comparison should be kept apart. The first is baseline comparison,
\[
d(X_t,X_0),
\]
where $X_0$ is a normative original state and $d$ measures historical deviation from it. This is the mode attributed to Rousseau. The second is internal self-relation. Let $\mathcal{G}(X_t)$ denote the capacities a configuration regenerates from its own operation, and $\mathcal{R}(X_t)$ the depletion of capacity attributable to its valuations. The net rate of change of the configuration's generative capacity $\mathcal{C}$ is then
\begin{equation}\label{eq:delta}
\Delta(X_t) \;=\; \frac{d\,\mathcal{C}}{dt} \;=\; \mathcal{G}(X_t) - \mathcal{R}(X_t).
\end{equation}
The depletion term has three components,
\[
\mathcal{R} \;=\; \mathcal{R}_{\mathrm{loss}} + \mathcal{R}_{\mathrm{inhibition}} + \mathcal{R}_{\mathrm{foregone}},
\]
corresponding to the destruction of an existing capacity, the prevention of its exercise, and the prevention of its formation. The distinction matters because the cases of interest do not all involve repression in the narrow sense.
A configuration is decadent when $\Delta(X_t)$ is persistently negative, that is, when its organization increasingly acts against the conditions of its own production. The definition is a rate and not a level, so it accommodates the case Section 3 allows: heavy constraint and sublimation, hence large $\mathcal{R}$, is compatible with a healthy configuration provided $\mathcal{G}$ keeps pace. What marks decadence is not the amount of repression but the failure of regeneration.

The third mode is comparison without privilege: ranking $X_t$ against $X_{t'}$ by their capacities or trajectories while neither serves as a norm. The transformational account requires this mode, since without it Nietzsche could pronounce a configuration decadent but never another one higher, and his positive evaluations of the Renaissance, of Napoleon, and of Goethe would have no footing.\footnote{\emph{Twilight of the Idols}, ``Skirmishes of an Untimely Man,'' \S48 (Napoleon) and \S49 (Goethe and the naturalness of the Renaissance).} What the account removes is the anchor and not the possibility of ranking. Comparison remains available so long as no term is privileged by priority.

Decadence has a second marker beyond the sign of $\Delta$. Let $\Gamma(X_t)$ denote the set of configurations reachable from $X_t$ under its own operation. A configuration can retain large present capacity while the set of futures open to it shrinks,
\[
\Gamma(X_{t+\tau}) \subsetneq \Gamma(X_t),
\]
and contraction of this kind is a further way in which a configuration acts against the conditions of its own production. It is also the sense in which some configurations open possibilities and others close them.

This also clarifies Nietzsche's hostility to what is better called the Rousseauian idiom than to Rousseau's own intentions. The target is a structure of inference that the doctrine invites, whoever draws it: Rousseau, his readers, or later movements that inherited his vocabulary. The objection is not merely that the idiom is mistaken about prehistoric humanity. The more consequential objection is that suffering is permitted to perform too much inferential work. Present injury becomes evidence of social corruption; social corruption implies a prior or underlying natural goodness; natural goodness supplies a standard of judgment; and the interests of those who suffer can thereby appear as the requirements of humanity itself. Nietzsche's genealogical interruption occurs between every pair of terms. Suffering is real without being self-interpreting. Oppression can be real without proving an original innocence. A grievance can be justified without thereby becoming a theory of human nature.

Voltaire consequently matters less as the possessor of a competing doctrine than as the representative of a competing intellectual operation. His skepticism can destroy an illusion without requiring the discovery of innocence underneath it. Nietzsche radicalizes that operation until even the comforting assumptions of the Enlightenment become available for dissolution. Neither providence, nature, equality, progress, nor history guarantees the direction of transformation.

Yet removing the origin creates a problem Nietzsche cannot evade. Once $X_0$ ceases to be normative, the judgment that $X_t$ is decadent requires another criterion. Nietzsche can no longer say simply that modernity has fallen. He must specify what is being lost, suppressed, exhausted, or rendered impossible, and why those capacities deserve the value he assigns them. Genealogy can expose the ancestry of a valuation, but ancestry alone cannot supply the valuation that replaces it.

The reflexive consequence is severe. Nietzsche's own diagnosis must enter the machinery he applies to Rousseau:
\[
\text{Nietzschean diagnosis}
\rightarrow
\text{conditions of its appeal}
\rightarrow
\text{interests expressed by it}
\rightarrow
\text{valuation embedded within it}.
\]
If Rousseau's philosophy may be read as a symptom of those for whom the existing order has become intolerable, Nietzsche's philosophy may equally be examined as a symptom of those for whom egalitarian modernity, domestication, or the contraction of rank has become intolerable. This symmetry does not refute either thinker. It is what consistency requires once philosophy has learned to ask what kind of life needs a proposition to be true.

The deepest disagreement between Rousseau and Nietzsche therefore concerns neither whether societies deteriorate nor whether suffering should matter. It concerns what can legitimately be inferred from transformation. The Rousseauian idiom, as Nietzsche receives it, makes loss intelligible through an origin. Nietzsche attempts the more difficult construction: loss without an original innocence, sickness without a prior state of grace, and decline without return.

For that reason revaluation, rather than restoration, is the necessary terminus of Nietzsche's account. Restoration moves backward along a presumed path toward a privileged state:
\[
X_t \rightarrow X_{t-1} \rightarrow \cdots \rightarrow X_0.
\]
Revaluation does not invert the trajectory. It introduces another transformation,
\[
X_t \xrightarrow{\ R\ } X_{t+1},
\]
whose legitimacy cannot come from resemblance to what preceded it. The future is not justified by recovering the past. It must produce and sustain capacities by criteria that remain themselves available to genealogical examination.

Nietzsche's correction of the language of a ``return to nature'' is therefore more than a rhetorical flourish. One does not go back. One climbs, transforms, experiments, and risks producing something for which history supplies no conserved template. The rejection of Rousseau's lost innocence leaves Nietzsche with a harsher conception of historical possibility: there was no state of grace to lose, and there is none to which we can return. If a healthier configuration is possible, it has to be made.

\appendix
\section{The Criterion Applied}

Nietzsche wrote about none of the institutions discussed below, so what follows is an application of the internal criterion of decadence, offered as a test of the apparatus and not as exegesis. The claim in each case is conditional: if $\Delta(X_t)$ is the right measure, then the following arrangements can be assessed by it.

\paragraph{Protection and load.} Nietzsche's maxim that what does not destroy a person makes that person stronger\footnote{\emph{Twilight of the Idols}, ``Maxims and Arrows,'' \S8.} has a structural resemblance to what is now called antifragility, the observation that some systems gain capacity from stressors and degrade in their absence. In the vocabulary of Section~3, capacities are regenerated under load, so an arrangement optimized to remove discomfort lowers $\mathcal{G}$ without any agent needing to repress anything. The valuation ``comfort'' consumes its own conditions. This is the point at which the charge of coddling acquires content: it names a configuration whose dominant valuation removes the stressors on which regeneration depends, and it is a claim about $\Delta$, not about softness of character.

\paragraph{Scaffold and cocoon.} Early childhood raises a harder case. Children are now offered choices continually, about food, clothing, and activities, which can foster the impression that the world is organized around them. That impression is false, but the arrangement producing it is not simply a defect. A protected and simplified environment lets a child isolate salient features quickly, and the simplification is productive precisely because it is a restriction. The transformational account has room for this, since it never treats constraint as bad in itself. The distinction that matters is responsiveness. Let $p(t)$ denote the degree of protection, $h(t)$ the environmental hazard, and $\mathcal{C}(t)$ the relevant capacity of the protected party. An appropriate protection level $p^{*}=f(\mathcal{C},h)$ satisfies
\[
\frac{\partial p^{*}}{\partial \mathcal{C}}<0,
\]
so that, holding hazard fixed, appropriate protection falls as capacity rises. Protection need not fall continuously. It can hold steady, rise when a new hazard appears, or fall in steps. An arrangement is scaffolding when $p$ tracks $p^{*}$, and a cocoon when
\[
p(t) > p^{*}\big(\mathcal{C}(t),h(t)\big)
\]
persistently, especially when the excess itself suppresses the development of $\mathcal{C}$. The cocoon contains a feedback loop:
\[
p\uparrow \;\Rightarrow\; \text{exposure}\downarrow \;\Rightarrow\; \dot{\mathcal{C}}\downarrow \;\Rightarrow\; \text{apparent need for } p\uparrow.
\]
The arrangement generates evidence for its own continuation, and that evidence is partly a product of the arrangement. This is endogenous self-undermining of the kind that defines decadence on the present account. The formalization is capacity-relative, since some capacities, language acquisition among them, depend on early exposure whatever the level of protection, and for those the relevant variable is not $p$ as defined here. Which description fits current upbringing is an empirical question this essay does not settle.

\paragraph{Pleasure as an end.} Nietzsche treats the pursuit of comfort and pleasure as a goal with sustained suspicion. The last man of the prologue to \emph{Thus Spoke Zarathustra} keeps a small pleasure for the day and another for the night,\footnote{\emph{Thus Spoke Zarathustra}, Prologue \S5.} and a maxim in \emph{Twilight of the Idols} says that only the Englishman strives for happiness.\footnote{\emph{Twilight of the Idols}, ``Maxims and Arrows,'' \S12; some translations render the German as pleasure and others as happiness.} Applied to contemporary tourism, theme parks, film, music, video games, and drugs, this suggests a diagnostic reading, though one that has to be stated carefully. Nietzsche praised music, dance, and laughter, and he distinguished intoxication that heightens capacity from narcotics that dull it. The criterion therefore cannot be the presence of pleasure. It is function. Recreation that returns the person to demanding work with capacities restored contributes to $\mathcal{G}$. Pleasure that substitutes for the tension from which capacities develop, and that must be renewed at increasing intensity, contributes to $\mathcal{R}$. The same activity can fall on either side depending on what it displaces.

\paragraph{Religion as supply of order.} Nietzsche's treatment of religion includes the view that it is a technique for relieving suffering, and the third essay of the \emph{Genealogy} describes the ascetic ideal as a means of deadening pain through intense feeling.\footnote{\emph{On the Genealogy of Morality}, Third Essay, \S15 (the priest as diverter of the course of ressentiment, since the sufferer seeks a responsible doer), \S16 (guilt and sin as the ideas through which ressentiment is turned backward), \S17 (relief of suffering and not treatment of its cause), and \S\S19--20 (emotional excess as a remedy for pain, applied under a religious interpretation). Nietzsche's own retrospective summary of the three essays is in \emph{Ecce Homo}, ``Books,'' on \emph{The Genealogy of Morals}.} He also treats faith as the satisfaction of a need for support and certainty. Extending this to a dependence on order and consolation that resembles addiction is a step beyond his text, but it follows the same logic: an arrangement is assessed by whether the relief it provides costs the person the capacities that suffering would otherwise have produced. Nietzsche does not reduce religion to this function. He credits the priestly and ascetic traditions with turning cruelty inward and thereby giving the human being an interior depth, a psychological complexity that the earlier condition lacked.\footnote{On the internalization of instinct and the acquisition of an inner world, \emph{On the Genealogy of Morality}, Second Essay, \S\S16--18; on the priestly form of existence as the soil on which the human became an interesting animal, First Essay, \S6.} The addiction reading captures the cost side of an assessment whose credit side is that sublimation.

\paragraph{The reflexive check.} The symptom method applies to these diagnoses too. Contempt for mass leisure is easily read as an expression of the position of those who dislike its tastes, and the vocabulary of decadence can serve as a rank marker, the same operation Nietzsche attributes to Rousseau with the direction reversed. By the comparative constraint of Section~3, the accusation is unfalsifiable unless $\Delta$ is actually assessed. An assessor must identify $\mathcal{C}$, specify the timescale, and name the counterfactual or comparator. A usable test is comparative and dispositional: whether a person's capacity for sustained difficulty is higher or lower after a period of a given consumption pattern than before it. Absent that comparison, the language of decadence describes a mood and not a configuration.

\section*{References}

\paragraph{Primary sources.} Nietzsche, \emph{Human, All Too Human} (first edition, 1878); \emph{Thus Spoke Zarathustra}; \emph{Beyond Good and Evil}; \emph{On the Genealogy of Morality}; \emph{Twilight of the Idols}; \emph{Ecce Homo}. Rousseau, \emph{Discourse on the Origin and Foundations of Inequality Among Mankind}; \emph{Of the Social Contract}. Nietzsche is cited by work, essay or part, and section number, which survive changes of translation and edition better than page references.

\paragraph{Secondary literature.}
\begin{itemize}
\item Janaway, Christopher. \emph{Beyond Selflessness: Reading Nietzsche's Genealogy}. Oxford University Press, 2007.
\item Kaufmann, Walter. \emph{Nietzsche: Philosopher, Psychologist, Antichrist}. Princeton University Press, 1950.
\item Leiter, Brian. \emph{Nietzsche on Morality}. Routledge, 2002.
\item Neuhouser, Frederick. \emph{Rousseau's Theodicy of Self-Love: Evil, Rationality, and the Drive for Recognition}. Oxford University Press, 2008.
\item Reginster, Bernard. \emph{The Affirmation of Life: Nietzsche on Overcoming Nihilism}. Harvard University Press, 2006.
\end{itemize}

\end{document}
